% Printed 498; principal PDF 505; Google PDF 509.
\par\clearpage
\noindent\makebox[0pt][l]{\fontsize{13}{16}\selectfont 498}\hfill{\small\textbf{BRIEFWECHSEL.}}\hfill\null\par\medskip
\noindent dass ich in meiner Sternwarte für solche Probemessungen keine recht geeignete Aufstellung finden kann, dass auch nur wenige geeignete Objekte sichtbar sind und dass schließlich die Winterzeit für meinen sehr empfindlichen Körper ungünstig ist. Ich will aber wenigstens versuchen, was in meinen Kräften steht.

\par\noindent\hspace*{\parindent}Sollten Sie vielleicht selbst ähnliche Erfahrungen gemacht und aus diesem Anlass verschiedene einschlägige Versuche angestellt haben, so würden Sie mich durch eine Mitteilung sehr verpflichten. Es ist tatsächlich eine Lebensfrage bei Theodoliten dieser Bauart. Bei \textsc{Borda}schen Kreisen kann der Fehler aufgrund des Konstruktionsprinzips gar nicht auftreten. Es gibt keinen Grund dafür, dass die Bewegung der Alhidade die Lage des unteren Fernrohrs \textls[120]{gegenüber dem Limbus}kreis verschiebt, und ebenso wenig verschiebt die Bewegung des ganzen Kreises die Lage der Alhidade gegenüber dem Limbuskreis. Man kann Theodoliten nach demselben Prinzip bauen, doch sind es dann ganz andere Instrumente als die \textsc{Reichenbach}schen. Wenn ich mich nicht irre, hat \textsc{Schumacher} einen solchen von \textsc{Gambey} (bei dem das untere Fernrohr dann eine ganz andere Rolle spielt). Ich kenne aber die Einrichtung im Einzelnen nicht.\enspace .\ .\ .\ .\ .\ .

\par\noindent\hspace*{\parindent}Es ist auffallend, wie Dinge zu Papier gebracht einen anderen Eindruck machen, als wenn man sie nur im Kopf durchdenkt. Wenn ich die obigen Zeilen noch einmal flüchtig überfliege, scheint mir, dass einem unbefangenen Leser, der selbst keinerlei eigene Erfahrungen gemacht hat, ein Umstand wichtiger erscheinen muss, als ich ihn bisher eingeschätzt habe. Es ist der Umstand, dass alle Theodoliten \textls[120]{zu Anfang} diese Erscheinung gar nicht oder nicht merklich gezeigt haben. Irgendeine \textls[120]{Abnutzung} muss also im Spiel sein, und ich werde deshalb wohl doch jedenfalls neue Schrauben, Muttern und Kugeln anfertigen lassen, obwohl sich, wie gesagt, an diesen Teilen einzeln kein Spiel erkennen lässt. Ich will aber versuchen, ob sich mit einem \textls[120]{viel stärkeren} Fernrohr nicht vielleicht schon an einem einzelnen Teil ein solches Spiel erkennen lässt.

\par\bigskip
\begin{center}{\large\textsc{Gauss} an \textsc{Bessel}.\quad Göttingen, 15. August 1844.}\end{center}
\par\noindent\hspace*{\parindent}.\ .\ .\ .\ .\ .\enspace Der Gegenstand, über den ich zuletzt mit Ihnen korrespondiert habe, nämlich die Neigung der Repetitions\textls[120]{theodoliten}, zu kleine Winkel anzu-

% Printed 499; principal PDF 506; Google PDF 510.
\par\clearpage
\noindent\null\hfill{\small\textbf{MESSUNGSFEHLER.}}\hfill\makebox[0pt][r]{\fontsize{13}{16}\selectfont 499}\par\medskip
\noindent geben, hat mich im vergangenen Winter noch sehr geplagt. Die Ursache kann keine andere sein, als dass der Kreis während derjenigen Bewegungen der Alhidade, bei denen er als fest vorausgesetzt wird, nicht absolut fest bleibt. Die Änderungen aber, die ich habe vornehmen lassen, haben sich als sehr wirksam erwiesen, sowohl bei den äußerst zahlreichen Probemessungen, die ich selbst von Januar bis März vorgenommen habe, als auch beim tatsächlichen Gebrauch, den mein Sohn weiterhin von diesem Instrument macht. Die Ergebnisse des Vorjahres waren zum Teil so schlecht, dass sie ganz verworfen werden mussten, während sie dieses Jahr sämtlich so gut sind, wie man es von einem Instrument dieser Größe nur erwarten kann. Die wichtigste Änderung besteht darin, dass der Alhidadenzapfen und seine Büchse (die ihrerseits den Limbuskreiszapfen bildet) jeweils von einer Tragfeder gestützt werden, sodass sich jede dieser Tragfedern unabhängig von der anderen nach Belieben und sofort spannen oder entspannen lässt (ganz oder teilweise). Im Gebrauch findet dann ein beständiger \textls[120]{planmäßiger} Wechsel zwischen den Zuständen dieser Federn statt, sodass insgesamt für jede Wiederholung 19 Handgriffe nötig sind (die Ablesung mitgerechnet). Man gewöhnt sich die Reihenfolge dieser Handgriffe jedoch so an, dass man fast ebenso schnell arbeiten kann wie mit der gewöhnlichen Einrichtung.
\par\bigskip
\begin{center}\rule{2.6cm}{.4pt}\end{center}
\vspace{1.2cm}
{\small
\begin{center}BEMERKUNGEN.\end{center}
\par\noindent\hspace*{\parindent}Die Notiz über die bei der hannoverschen Gradmessung und Landesvermessung benutzten und die für die letztere noch erforderlichen Instrumente ist einem im Gauss-Archiv befindlichen Bericht von \textsc{Gauss} vom 21. März 1829 an den Geheimen Kabinettsrat \textsc{Hoppenstedt} entnommen.
\par\noindent\hspace*{\parindent}Zu den vorstehenden Briefen an \textsc{Olbers}, \textsc{Schumacher} und \textsc{Bessel}, die nach den Originalen abgedruckt sind, sind die Stationsausgleichung für Brillit, S. 265/267, sowie eine Bemerkung von Professor \textsc{Goldschmidt}, S. 289, zu vergleichen.
\par\noindent\hfill\textsc{Krüger.}\hspace*{1em}\par}
\vspace{1cm}
\begin{center}\rule{3.5cm}{.4pt}\end{center}
\par\medskip\noindent\makebox[.86\linewidth][r]{63\textsuperscript{*}}\par

% Printed 500; principal PDF 507; Google PDF 511.
\par\clearpage
\noindent\makebox[0pt][l]{\fontsize{13}{16}\selectfont 500}\hfill{\small\textbf{BRIEFWECHSEL.}}\hfill\null\par
\vspace*{3.2cm}
\begin{center}
{\LARGE BRIEFWECHSEL UND NACHLASS.}\par
\vspace{1.5cm}
{\large [Berechnung des mittleren Ablesefehlers einschließlich des Teilungsfehlers\par am Theodoliten.]}\par
\bigskip\bigskip
{\large [1.]}\par\medskip
{\large\textsc{Gauss} an \textsc{Gerling}.\quad Göttingen, 17. April 1844.}
\end{center}
\par\noindent\hspace*{\parindent}.\ .\ .\ .\ .\ .\enspace Über eine andere Kontrollrechnung an Theodolitenbeobachtungen (die den mittleren Fehler der \textls[120]{Ablesungen} einschließlich des Teilungsfehlers bestimmen soll) muss ich Ihnen ein anderes Mal schreiben. Ich habe diese Rechnungen schon vor 19 Jahren ausgeführt, und sie haben mir viel Vergnügen gemacht. Ich habe nämlich für \textls[120]{jede} vorgekommene vollständige Ablesung $a-b+c-d=l$ berechnet, wobei (unter Weglassung der Grade) $a$, $b$, $c$, $d$ die 4 Ablesungen bezeichnen. Dieses $l$ sollte (\textls[120]{trotz} der Exzentrizität, ja trotz einer möglicherweise veränderlichen Exzentrizität) \textls[120]{konstant} sein. Es ist nur infolge der Ablese- und Teilungsfehler veränderlich. Setzt man den Mittelwert aus sehr vielen, $n$, Ablesungen $=\lambda$ (der sich bei einem gegebenen Theodoliten aus vielen Beobachtungen sehr genau bestimmen lässt und unveränderlich sein sollte) und $\frac{\Sigma(l-\lambda)^2}{n-1}=kk$, so ist $\frac12k$ der mittlere Fehler, mit dem man bei der Ablesung eines Index, und $\frac14k$ der mittlere Fehler, mit dem man beim Mittelwert aus allen 4 Indizes rechnen muss. Ich finde nun den letzteren oder $\frac14k$ aus vollständigen Ablesungen:

% Printed 501; principal PDF 508; Google PDF 512.
\par\clearpage
\noindent\null\hfill{\small\textbf{MESSUNGSFEHLER.}}\hfill\makebox[0pt][r]{\fontsize{13}{16}\selectfont 501}\par\medskip
\begin{center}
\begingroup\renewcommand{\arraystretch}{1.3}\setlength{\tabcolsep}{11pt}
\begin{tabular}{@{\vrule width .8pt\hspace{11pt}}c|c|c@{\hspace{11pt}\vrule width .8pt}}
\noalign{\hrule height .8pt}
[Im Jahr & Anzahl der Ables. & Mittl. Fehler]\\[4pt]
\hline
1821 & \makebox[2.4em][r]{410} & \makebox[3em][r]{$1\overset{\prime\prime}{,}25$}\\[7pt]
1822 & \makebox[2.4em][r]{717} & \makebox[3em][r]{$1{,}57$}\\
1823 & \makebox[2.4em][r]{449} & \makebox[3em][r]{$1{,}61$}\\
1824 & \makebox[2.4em][r]{843} & \makebox[3em][r]{$1{,}60$}\\
1825 & \makebox[2.4em][r]{609} & \makebox[3em][r]{$1{,}55$}\\
\hline
1822---1825 & \makebox[2.4em][r]{2618} & \makebox[3em][r]{$1{,}58$}\\
\end{tabular}
\endgroup
\end{center}
\par\noindent\hspace*{\parindent}Aus der so engen Übereinstimmung der Jahre 1822---1825 lässt sich bereits schließen, dass der abweichende Wert von 1821 nicht zufällig ist. Tatsächlich waren die Teilstriche an \textsc{Schumacher}s Theodoliten noch etwas schöner als an meinem eigenen. Jener stammte noch aus der Zeit, als \textsc{Reichenbach} selbst die Werkstatt hatte. Meiner war bei \textsc{Reichenbach} bestellt, wurde aber von \textsc{Ertel} unter seiner Aufsicht angefertigt. Auch an meinem 8-zölligen Theodoliten (von 1829) ist die Teilung merklich weniger schön als an dem 8-zölligen von \textsc{Reichenbach}, den ich 1813 erworben habe. Jener hat aber ein \textls[120]{viel} besseres .Fernrohr.

\par\noindent\hspace*{\parindent}Ich sehe, dass ich Ihnen nun doch das Wesentlichste über diesen Gegenstand geschrieben habe, und erwähne daher nur noch der Vollständigkeit halber eine wesentliche Abkürzung jener Rechnung. Es ist nämlich $\Sigma l=n\lambda$, also
\[
\Sigma(l-\lambda)^2=\Sigma ll-2n\lambda\lambda+n\lambda\lambda=\Sigma ll-n\lambda\lambda
\left[=\Sigma ll-\frac{(\Sigma l)^2}{n}\right],
\]
\noindent also
\[
k=\frac{\Sigma ll}{n-1}-\frac{n}{n-1}\lambda\lambda
\left[=\frac{\Sigma ll}{n-1}-\frac{(\Sigma l)^2}{n(n-1)}\right].
\]
\par\bigskip\bigskip
\begin{center}{\large [2.]}\end{center}
\par\noindent\hspace*{\parindent}Zur Ermittlung der Fehler, die sich aus der Teilung und dem Ablesen ergeben, wurden alle vollständigen Ablesungen am Theodoliten, 449 an der Zahl [aus dem Jahr 1823], auf folgende Weise ausgewertet.

% Printed 502; principal PDF 509; Google PDF 513.
\par\clearpage
\noindent\makebox[0pt][l]{\fontsize{13}{16}\selectfont 502}\hfill{\small\textbf{NACHLASS.}}\hfill\null\par\medskip
\par\noindent\hspace*{\parindent}Summe der Ablesungen am ersten und dritten Index minus Summe der Ablesungen am zweiten und vierten $=l$.
\[
\begin{aligned}
\Sigma l&=+133\\
\Sigma ll&=\phantom{+}18521,
\end{aligned}
\]
\noindent also der Mittelwert von $l=\left[\frac{133}{449}=\right]+0\overset{\prime\prime}{,}30=\lambda$;
\[
\text{Mittelwert von }(l-\lambda)^2=
\left[\frac{1}{448}\left(18521-\frac{133^2}{449}\right)=\right]41{,}254.
\]
\par\noindent\hspace*{\parindent}Also der mittlere Fehler bei einem Nonius einschließlich des Teilungsfehlers:
\[
\left[\frac{\sqrt{41{,}254}}{2}=\right]\pm3\overset{\prime\prime}{,}211.
\]
\par\noindent\hspace*{\parindent}An \textsc{Schumacher}s Theodoliten ergaben 410 Ablesungen aus dem Jahr 1821
\[
\begin{aligned}
\Sigma l&=-1896\\
\Sigma ll&=\phantom{-}18982.
\end{aligned}
\]
\par\noindent\hspace*{\parindent}Mittelwert von $l=-4\overset{\prime\prime}{,}62=\lambda$.
\par\noindent\hspace*{\parindent}Mittelwert von $(l-\lambda)^2=24{,}97$.
\par\noindent\hspace*{\parindent}Mittlerer Ablesefehler eines Nonius einschließlich des Teilungsfehlers $=\pm2\overset{\prime\prime}{,}499$.
\par\bigskip\bigskip
\par\noindent\hspace*{\parindent}Der Theodolit ergab 1822 aus 717 Ablesungen
\[
\begin{aligned}
\Sigma l&=+212\\
\Sigma ll&=\phantom{+}28433.
\end{aligned}
\]
\par\noindent\hspace*{\parindent}Mittelwert von $l=+0\overset{\prime\prime}{,}30=\lambda$.
\par\noindent\hspace*{\parindent}Mittelwert von $(l-\lambda)^2=39{,}623$.
\par\noindent\hspace*{\parindent}Mittlerer Ablesefehler bei einem Nonius einschließlich des Teilungsfehlers $=\pm3\overset{\prime\prime}{,}147$.
\par\bigskip\bigskip
\par\noindent\hspace*{\parindent}Dieses Ergebnis stimmt so eng mit dem von 1823 überein, dass wir beide zusammenfassen können. Wir erhalten also aus 1166 Ablesungen:

% Printed 503; principal PDF 510; Google PDF 514.
\par\clearpage
\noindent\null\hfill{\small\textbf{MESSUNGSFEHLER.}}\hfill\makebox[0pt][r]{\fontsize{13}{16}\selectfont 503}\par\medskip
\[
\begin{aligned}
\Sigma l&=+345\\
\Sigma ll&=\phantom{+}46954\\
\lambda&=+0\overset{\prime\prime}{,}296.
\end{aligned}
\]
\par\noindent\hspace*{\parindent}Mittelwert von $(l-\lambda)^2=40{,}216$.
\par\noindent\hspace*{\parindent}Mittlerer Fehler bei einem Nonius $=\pm3\overset{\prime\prime}{,}1708$.
\par\noindent\hspace*{\parindent}Die Häufigkeiten der einzelnen Werte von $l$ ohne Rücksicht auf das Vorzeichen waren:
\begin{center}
\begingroup\renewcommand{\arraystretch}{1.15}\setlength{\tabcolsep}{10pt}
\begin{tabular}{@{\vrule width .8pt\hspace{10pt}}c|c@{\hspace{10pt}\vrule width .8pt\hspace{10pt}}c|c@{\hspace{10pt}\vrule width .8pt\hspace{10pt}}c|c@{\hspace{10pt}\vrule width .8pt}}
\noalign{\hrule height .8pt}
\makebox[0pt][l]{$[$}\makebox[2em]{$l$} & Anzahl & \makebox[2em]{$l$} & Anzahl & \makebox[2em]{$l$} & Anzahl $]$\\[4pt]
\hline
\makebox[.6em][r]{0}\rlap{$''$} & \makebox[1.8em][r]{45} & \makebox[1.2em][r]{6}\rlap{$''$} & \makebox[1.8em][r]{133} & \makebox[1.2em][r]{11}\rlap{$''$} & \makebox[1.2em][r]{38}\\
\makebox[.6em][r]{1} & \makebox[1.8em][r]{106} & \makebox[1.2em][r]{7} & \makebox[1.8em][r]{104} & \makebox[1.2em][r]{12} & \makebox[1.2em][r]{24}\\
\makebox[.6em][r]{2} & \makebox[1.8em][r]{108} & \makebox[1.2em][r]{8} & \makebox[1.8em][r]{82} & \makebox[1.2em][r]{13} & \makebox[1.2em][r]{15}\\
\makebox[.6em][r]{3} & \makebox[1.8em][r]{123} & \makebox[1.2em][r]{9} & \makebox[1.8em][r]{84} & \makebox[1.2em][r]{14} & \makebox[1.2em][r]{9}\\
\makebox[.6em][r]{4} & \makebox[1.8em][r]{113} & \makebox[1.2em][r]{10} & \makebox[1.8em][r]{52} & \makebox[1.2em][r]{15} & \makebox[1.2em][r]{4}\\
\makebox[.6em][r]{5} & \makebox[1.8em][r]{126} &  &  &  & \\
\end{tabular}
\endgroup
\end{center}
\par\noindent\hspace*{\parindent}Im Jahr 1824 ergab der Theodolit aus 843 Ablesungen
\[
\begin{aligned}
\Sigma l&=+104\\
\Sigma ll&=\phantom{+}34506\\
\lambda&=+0\overset{\prime\prime}{,}12.
\end{aligned}
\]
\par\noindent\hspace*{\parindent}Mittelwert von $(l-\lambda)^2=40{,}966$.
\par\noindent\hspace*{\parindent}Mittlerer Fehler bei einem Nonius einschließlich des Teilungsfehlers $=\pm3\overset{\prime\prime}{,}200$.
\par\noindent\hspace*{\parindent}Also aus allen 3 Jahren [1822---1824] mit 2009 Ablesungen
\[
\begin{aligned}
\Sigma l&=+449\\
\Sigma ll&=\phantom{+}81460\\
\lambda&=+0\overset{\prime\prime}{,}2235.
\end{aligned}
\]
\par\noindent\hspace*{\parindent}Mittelwert von $(l-\lambda)^2=40{,}518$.
\par\noindent\hspace*{\parindent}Mittlerer Fehler eines Nonius $=\pm3\overset{\prime\prime}{,}1827$. Die Häufigkeiten der einzelnen Werte von $l$ ohne Rücksicht auf das Vorzeichen waren:

% Printed 504; principal PDF 511; Google PDF 515.
\par\clearpage
\noindent\makebox[0pt][l]{\fontsize{13}{16}\selectfont 504}\hfill{\small\textbf{NACHLASS.}}\hfill\null\par\medskip
\begin{center}
\begingroup\renewcommand{\arraystretch}{1.15}\setlength{\tabcolsep}{10pt}
\begin{tabular}{@{\vrule width .8pt\hspace{10pt}}c|c@{\hspace{10pt}\vrule width .8pt\hspace{10pt}}c|c@{\hspace{10pt}\vrule width .8pt\hspace{10pt}}c|c@{\hspace{10pt}\vrule width .8pt}}
\noalign{\hrule height .8pt}
\makebox[0pt][l]{$[$}\makebox[2em]{$l$} & Anzahl & \makebox[2em]{$l$} & Anzahl & \makebox[2em]{$l$} & Anzahl $]$\\[4pt]
\hline
\makebox[.6em][r]{0}\rlap{$''$} & \makebox[1.8em][r]{80} & \makebox[1.2em][r]{7}\rlap{$''$} & \makebox[1.8em][r]{191} & \makebox[1.2em][r]{14}\rlap{$''$} & \makebox[1.2em][r]{13}\\
\makebox[.6em][r]{1} & \makebox[1.8em][r]{182} & \makebox[1.2em][r]{8} & \makebox[1.8em][r]{155} & \makebox[1.2em][r]{15} & \makebox[1.2em][r]{6}\\
\makebox[.6em][r]{2} & \makebox[1.8em][r]{185} & \makebox[1.2em][r]{9} & \makebox[1.8em][r]{138} & \makebox[1.2em][r]{16} & \makebox[1.2em][r]{1}\\
\makebox[.6em][r]{3} & \makebox[1.8em][r]{195} & \makebox[1.2em][r]{10} & \makebox[1.8em][r]{101} & \makebox[1.2em][r]{17} & \makebox[1.2em][r]{1}\\
\makebox[.6em][r]{4} & \makebox[1.8em][r]{205} & \makebox[1.2em][r]{11} & \makebox[1.8em][r]{63} & \makebox[1.2em][r]{18} & \makebox[1.2em][r]{---}\\
\makebox[.6em][r]{5} & \makebox[1.8em][r]{217} & \makebox[1.2em][r]{12} & \makebox[1.8em][r]{38} & \makebox[1.2em][r]{19} & \makebox[1.2em][r]{1}\\
\makebox[.6em][r]{6} & \makebox[1.8em][r]{214} & \makebox[1.2em][r]{13} & \makebox[1.8em][r]{23} &  & \\
\end{tabular}
\endgroup
\end{center}
\begin{center}\rule{3.5cm}{.4pt}\end{center}
\par\bigskip
\par\noindent\hspace*{\parindent}Am [\textsc{Borda}schen] Kreis wurden [1823] 46 vollständige Ablesungen vorgenommen. Es ergab sich
\[
\begin{aligned}
\Sigma l&=-174\\
\Sigma ll&=\phantom{-}1534\\
\lambda&=-3\overset{\prime\prime}{,}78.
\end{aligned}
\]
\par\noindent\hspace*{\parindent}Mittelwert von $(l-\lambda)^2=19{,}46$.
\par\noindent\hspace*{\parindent}Mittlerer Ablesefehler eines Nonius einschließlich des Teilungsfehlers $=\pm2\overset{\prime\prime}{,}206$.
\par\noindent\hspace*{\parindent}Der Kreis ergab 1821 aus 92 Ablesungen
\[
\begin{aligned}
\Sigma l&=-650\\
\Sigma ll&=\phantom{-}11578\\
\lambda&=-7\overset{\prime\prime}{,}07.
\end{aligned}
\]
\par\noindent\hspace*{\parindent}Mittelwert von $(l-\lambda)^2=76{,}767$.
\par\noindent\hspace*{\parindent}Mittlerer Ablesefehler eines Nonius einschließlich des Teilungsfehlers $=\pm4\overset{\prime\prime}{,}381$.
\par\noindent\hspace*{\parindent}Am Kreis wurden 1822 insgesamt 137 Ablesungen vorgenommen. Diese ergaben
\[
\begin{aligned}
\Sigma l&=-543\\
\Sigma ll&=\phantom{-}6155\\
\lambda&=-3\overset{\prime\prime}{,}96.
\end{aligned}
\]
\par\noindent\hspace*{\parindent}Mittelwert von $(l-\lambda)^2=29{,}432$.

% Printed 505; principal PDF 512; Google PDF 516.
\par\clearpage
\noindent\null\hfill{\small\textbf{MESSUNGSFEHLER.}}\hfill\makebox[0pt][r]{\fontsize{13}{16}\selectfont 505}\par\medskip
\par\noindent\hspace*{\parindent}Mittlerer Ablesefehler bei einem Nonius einschließlich des Teilungsfehlers $=\pm2\overset{\prime\prime}{,}713$.
\par\noindent\hspace*{\parindent}Fasst man die Messungen von 1822 und 1823 zusammen, so ergibt sich [bei 183 Ablesungen]
\[
\begin{aligned}
\Sigma l&=-717\\
\Sigma ll&=\phantom{-}7689\\
\lambda&=-3\overset{\prime\prime}{,}92.
\end{aligned}
\]
\par\noindent\hspace*{\parindent}Mittelwert von $(l-\lambda)^2=\left[26{,}812.\right.$
\par\noindent\hspace*{\parindent}Mittlerer Ablesefehler bei einem Nonius einschließlich des Teilungsfehlers $=\pm2\overset{\prime\prime}{,}589.]$
\par\bigskip\bigskip
\begin{center}\rule{2.6cm}{.4pt}\end{center}
\vspace{1.5cm}
{\small
\begin{center}BEMERKUNGEN.\end{center}
\par\noindent\hspace*{\parindent}Der vorstehende Brief an \textsc{Gerling}, Art. [1], wurde nach dem Original abgedruckt. Art. [2] wurde einem Beobachtungsheft für die hannoversche Gradmessung entnommen.
\par\noindent\hfill\textsc{Krüger.}\hspace*{1em}\par}
\vspace{2cm}
\begin{center}\rule{3.5cm}{.4pt}\end{center}

\vfill
\noindent IX.\hfill 64\hspace*{1cm}\par

% Contextual printed 506; principal PDF 513; Google PDF 517: unnumbered blank leaf.
\par\clearpage
\null

% Contextual printed 507; principal PDF 514; Google PDF 518: unnumbered section title.
\par\clearpage
\vspace*{4cm}
\begin{center}
{\fontsize{22}{30}\selectfont
REDUKTION SCHIEFER WINKEL\par
AUF DEN HORIZONT.\par
REPETITIONSBEOBACHTUNGEN.\par}
\end{center}
\vspace{8.2cm}
\noindent\makebox[.86\linewidth][r]{64\textsuperscript{*}}\par
